% ================================================================
\section{The Parseval Bridge as Unitarity}
\label{sec:parseval}
% ================================================================

The key technical innovation of the Cathedral is the \emph{Parseval Bridge},
which rewrites the position-space $L^2$ norm as a frequency-space integral:

\[
  \underbrace{\int_0^1 |1 - f_{N,v}(x)|^2\,dx}_{\text{position space (vacuum energy)}}
  = \frac{1}{2\pi}\underbrace{\int_{-\infty}^{\infty}
  \left|\frac{1}{1/2+it} - \sum_k \frac{v_k}{k^{1/2+it}(1/2+it)}\right|^2\,dt}_{\text{frequency space (critical line)}}
\]

\begin{principle}[Unitarity of the S-Matrix]
The Parseval equality is the \emph{optical theorem} of the prime number
quantum field. In particle physics, the optical theorem states
\[
  2\,\im\,\mathcal{A}(\theta=0) = \sigma_{\mathrm{tot}}
\]
relating the forward scattering amplitude to the total cross section.
Both are consequences of unitarity: probability is conserved.

The Parseval Bridge says that the ``probability'' (squared norm) in
position space \emph{exactly} equals the ``probability'' in frequency
space. This is not approximate---it is an identity. The frequencies
are the values on the critical line $\re s = 1/2$, which is the
``mass shell'' of the Riemann zeta function.
\end{principle}

\subsection{The Three-Term Decomposition as Feynman Diagrams}

The integrand on the critical line decomposes as:
\[
  |1 - \zeta(s)W_N(s)|^2 / |s|^2
  = \underbrace{1/|s|^2}_{\text{free propagator}}
  - \underbrace{2\re(\zeta W_N)/|s|^2}_{\text{interaction}}
  + \underbrace{|\zeta W_N|^2/|s|^2}_{\text{self-energy}}
\]

\begin{correspondence}[Feynman Diagram Decomposition]
\begin{itemize}
  \item \textbf{Term 1} ($1/|s|^2$): The \emph{free propagator} of the
    massless field. Evaluates to exactly $1$ via the Cauchy distribution
    integral $\frac{1}{2\pi}\int_{-\infty}^\infty \frac{dt}{1/4+t^2} = 1$.
    This is the \emph{Cauchy resonance}---a Breit-Wigner peak at $t=0$
    with width $\Gamma = 1/2$, corresponding to the decay rate of
    the zeta function's pole at $s=1$.

    \emph{Cathedral code:} \lean{term1\_exact} in \lean{PlancherelBypass.lean}.

  \item \textbf{Term 2} ($-2\re(\zeta W_N)/|s|^2$): The \emph{interference}
    between the free field and the scattered wave. This is the analogue of
    the Born approximation in scattering theory. It requires contour shifting
    from $\re s = 1/2$ to $\re s = 2$, crossing the pole at $s=1$---a
    \emph{resonance crossing} that produces the $\ln\ln N$ correction.

  \item \textbf{Term 3} ($|\zeta W_N|^2/|s|^2$): The \emph{self-energy}
    of the scattered wave. Bounded by the Montgomery--Vaughan mean value
    theorem, which is the number-theoretic analogue of the \emph{Dyson
    equation} for the dressed propagator.
\end{itemize}
\end{correspondence}

\subsection{The Mellin--Perron Bridge as Bohr Complementarity}

The Cathedral's dual-path architecture---Mellin (frequency domain) and
spatial/Perron (position domain)---independently proves $d_N^2 \to 0$.
The theorem \lean{critical\_line\_mellin\_variance\_from\_perron}
(\lean{MellinPerronBridge.lean}, proved April~27, 2026) makes their
equivalence explicit via the Parseval isometry.

\begin{principle}[Bohr Complementarity]
The Mellin--Perron Bridge realizes Bohr's complementarity principle:
the vacuum energy $d_N^2$ is a \emph{basis-independent observable},
measurable equally well in position coordinates
($\int_0^1 |r_N(x)|^2\,dx$) or momentum coordinates
($\frac{1}{2\pi}\int_{-\infty}^{\infty} |M_{r_N}(1/2+it)|^2\,dt$).
The Bridge is the formal proof that the prime number quantum field
theory is self-consistent across representations.

In quantum mechanics, complementarity asserts that position and
momentum measurements are related by the Fourier transform.
Here, the spatial path (Perron contour integration, Mertens bound,
Abel summation) and the spectral path (Mellin transform, critical-line
variance) are the position and momentum representations of the same
physical system, and the Parseval isometry guarantees they yield
identical vacuum energies.

The two paths correspond to different \emph{gauge choices}:
\begin{center}
\begin{tabular}{@{}lllc@{}}
\toprule
\textbf{Path} & \textbf{Domain} & \textbf{Gauge} & \textbf{Axioms} \\
\midrule
Mellin (\lean{\_mellin}) & Frequency & Unitary (compact) & 2 \\
Spatial (\lean{\_spatial}) & Position & Lorenz (transparent) & 4 \\
Bridge (\lean{MellinPerronBridge}) & --- & Gauge transformation & 0 \\
\bottomrule
\end{tabular}
\end{center}
The unitary gauge (Mellin) uses fewer variables but the axioms are
composite. The Lorenz gauge (spatial) uses more variables but each
axiom maps to a single classical theorem. The Bridge proves the
gauge transformation with zero axioms.
\end{principle}

\subsection{The Stained Glass Rotors: Spectral Energy Partition}

The Stained Glass Rotors (\lean{GallagherPartition.lean},
\lean{GallagherMVT.lean}, \lean{FrequencySeparation.lean}, all
proved with zero sorry) provide a character-based decomposition of the
critical-line spectral energy.

\begin{correspondence}[Quantum Error-Correcting Code]
The discrete energy partition (\lean{discrete\_energy\_partition}, proved):
\[
  \sum_n |a_n|^2 = \frac{1}{4}\sum_{i=1}^{4}
  \sum_{n} |\chi_i(n)|^2 \cdot |a_n|^2
\]
decomposes the total Dirichlet polynomial energy into four orthogonal
channels indexed by Dirichlet characters $\chi_i$ modulo~8. This is a
\emph{stabilizer code}: the characters act as syndrome measurements,
and the fourfold energy split asserts that the prime lattice distributes
its energy uniformly across all syndrome sectors---\emph{geometric
frustration} at the arithmetic level.

The orthogonality of the four channels ($\sum_n \chi_i(n)\overline{\chi_j(n)}
= 0$ for $i \neq j$) is verified by \lean{native\_decide}---the Lean kernel
performs an exhaustive computation over all $8 \times 8$ character values.
\end{correspondence}

\begin{correspondence}[Gallagher MVT as Spectral Isolation]
The Gallagher Mean Value Theorem (\lean{gallagher\_mvt}, proved):
\[
  \int_{-\infty}^{\infty} \left|\sum_n a_n\, e^{i\lambda_n t}\right|^2
  \cdot K_\delta(t)\,dt = \sum_n |a_n|^2
\]
where $K_\delta$ is the Fej\'er kernel with width $\delta$, states that
the energy of a trigonometric polynomial equals the sum of squared
amplitudes---the number-theoretic \emph{completeness relation}. The
Dirichlet modes form a quasi-orthogonal set on the critical line.
\end{correspondence}

\begin{correspondence}[Log-Frequency Separation as Dispersion Relation]
The frequencies $\lambda_n = -\log(n)/(2\pi)$ of the Dirichlet
polynomial form a non-uniform lattice. The separation bound
$|\lambda_i - \lambda_j| \geq 1/(N+1)$ for $i \neq j$
(\lean{log\_frequencies\_separated}, proved) is the
\emph{dispersion relation} of the prime lattice: the energy levels
$(\log n)$ grow logarithmically, ensuring the spectral resolution
$\delta = 1/(N+1)$ suffices for mode decomposition. This mirrors
the Van Hove singularity in solid-state physics---the density of
states peaks where the dispersion relation flattens.
\end{correspondence}

